\section{USAMO 1996}
        \begin{enumerate}
        	\item (\textit{USAMO 1996}) Prove that the average of the numbers $n \sin n^{\circ} \; (n = 2,4,6,\ldots,180)$ is $\cot 1^{\circ}$.


 



\textbf{Solution 1}

First, as $180\sin{180^\circ}=0,$ we omit that term. Now, we multiply by $\sin 1^\circ$ to get, after using product to sum, $(\cos 1^\circ-\cos 3^\circ)+2(\cos 3^\circ-\cos 5^\circ)+\cdots +89(\cos 177^\circ-\cos 179^\circ)$. This simplifies to $\cos 1^\circ+\cos 3^\circ +\cos 5^\circ+\cos 7^\circ+...+\cos 177^\circ-89\cos 179^\circ$. Since $\cos x=-\cos(180-x),$ this simplifies to $90\cos 1^\circ$. We multiplied by $\sin 1^\circ$ in the beginning, so we must divide by it now, and thus the sum is just $90\cot 1^\circ$, so the average is $\cot 1^\circ$, as desired.

 $\Box$

 



\textbf{Solution 2}

Notice that for every $n\sin n^\circ$ there exists a corresponding pair term $(180^\circ - n)\sin{180^\circ - n} = (180^\circ - n)\sin n^\circ$, for $n$ not $90^\circ$. Pairing gives the sum of all $n\sin n^\circ$ terms to be $90(\sin 2^\circ + \sin 4^\circ + ... + \sin 178^\circ)$, and thus the average is 
 \[S = (\sin 2^\circ + \sin 4^\circ + ... + \sin 178^\circ). (*)\] We need to show that $S = \cot 1^\circ$. Multiplying (*) by $2\sin 1^\circ$ and using sum-to-product and telescoping gives $2\sin 1^\circ S = \cos 1^\circ - \cos 179^\circ = 2\cos 1^\circ$. Thus, $S = \frac{\cos 1^\circ}{\sin 1^\circ} = \cot 1^\circ$, as desired.

 $\Box$

 



\textbf{Solution 3 -hashtagmath}

We know that the average of a list of numbers is the sum of all the terms, divided by the number of terms. So we will set up an average. This average will average to $\cot 1^\circ$. So we can set it equal to $\frac{\cos 1^\circ}{\sin 1^\circ}$. Doing so, gives us 
 \[\frac{2 \sin 2^\circ + 4 \sin 4^\circ + 6 \sin 6^\circ + 8 \sin 8^\circ + 10 \sin 10^\circ + 12 \sin 12^\circ ... + 180 \sin 180^\circ}{90} = \frac{\cos 1^\circ}{\sin 1^\circ}.(1)\]

 We should try to simplify the numerator as it looks pretty messy.

 Now we know that $\sin x = \sin (180-x)$, so maybe we can use this to clear things up a bit. Applying this to $(1)$ gives us 
 \[\frac{2 \sin 178^\circ + 4 \sin 176^\circ + 6 \sin 174^\circ + 8 \sin 172^\circ + 10 \sin 170^\circ + 12 \sin 168^\circ ... + 180 \sin 0^\circ}{90} = \frac{\cos 1^\circ}{\sin 1^\circ}. (2)\]

 Now, we know that $\sin 2^\circ = \sin 178^\circ$, and $\sin 4^\circ = \sin 176^\circ$ and so on. We also know that $\sin 0^\circ = 0$, so we can omit the last term. Also, we know that $\sin 90 = 1$. Thus $90 \sin 90 = 90$. So we will use all these when we further simplify. Another thing that we can simplify is to notice that when we reach the term $92\sin 88^\circ$, we can rewrite it as $92\sin 92^\circ$. Thus we can rewrite everything in the form $x\sin y^\circ$ as $x \sin x^\circ$, starting with the term $92\sin 88^\circ$, then $94\sin 86^\circ$, and so on. 

 Thus we can further simplify our equation. Doing so, gives us 
 \[\frac{2 \sin 178^\circ + 4 \sin 176^\circ + 6 \sin 174^\circ + 8 \sin 172^\circ + 10 \sin 170^\circ + 12 \sin 168^\circ ...+ 176 \sin 176^\circ + 178 \sin 178^\circ + 90}{90} = \frac{\cos 1^\circ}{\sin 1^\circ}. (3)\]

 Then, to make things easier, we can rewrite all the numbers with degrees less than $90^\circ$ and combine like terms. Doing so gives us  
 \[\frac{180 \sin 2^\circ + 180 \sin 4^\circ + 180 \sin 6^\circ +180 \sin 8^\circ +180 \sin 10^\circ ... +180 \sin 88^\circ + (90)}{90} = \frac{\cos 1^\circ}{\sin 1^\circ}. (4)\]

 Now we can factor out $180$ to get 
 \[\frac{180(\sin 2^\circ+ \sin 4^\circ + \sin 6^\circ + \sin 8^\circ + \sin 10^\circ ... + \sin 88^\circ) + 90}{90} = \frac{\cos 1^\circ}{\sin 1^\circ}. (5)\]

 Then we can simplify it by getting rid of the denominator to get to 
 \[2(\sin 2^\circ+ \sin 4^\circ + \sin 6^\circ + \sin 8^\circ + \sin 10^\circ ... + \sin 88^\circ) + 1 = \frac{\cos 1^\circ}{\sin 1^\circ}. (6)\]

 Now we notice that we can multiply both sides by $\sin 1^\circ$. Doing so gives us 
 \[\sin 1^\circ(2 (\sin 2^\circ+ \sin 4^\circ + \sin 6^\circ + \sin 8^\circ + \sin 10^\circ ... + \sin 88^\circ) + 1) = \cos 1^\circ. (7)\]

 Now, simplifying $\sin 2^\circ+ \sin 4^\circ + \sin 6^\circ + \sin 8^\circ + \sin 10^\circ ... + \sin 88^\circ$ doesn't look too promising. So maybe if we expand again, we can maybe somehow use our product to sum formulas. 

 Doing so, gives us 
 \[(\sin 1^\circ)(2 \sin 2^\circ)+(\sin 1^\circ)(2\sin 4^\circ) +(\sin 1^\circ)(2\sin 6^\circ) + (\sin 1^\circ)(2\sin 8^\circ) +(\sin 1^\circ)(2 \sin 10^\circ) ... + (\sin 1^\circ)(2\sin 88^\circ) + \sin 1^\circ = \cos 1^\circ. (8)\]

 Now we can use our product of sines to sum formula and see if we can find a pattern. 

 We will start with expanding $(\sin 1^\circ)(2 \sin 2^\circ)$. Because the format of our formula for the product of sines is $\sin \alpha + \sin \beta$, we can factor out the $2$ and find the product, then multiply by $2$. So, we are at $2(\sin 1^\circ)(\sin 2^\circ)$. Our formula for the product of sines is $\sin \alpha + \sin \beta = \frac{1}{2}[\cos(\alpha - \beta) - \cos(\alpha + \beta)]$. Plugging in our values into the formula and simplifying gives us $2\left(\frac{1}{2}(\cos 1^\circ - \cos 3^\circ)\right)$. We know that $2\left(\frac{1}{2}\right)$ will just cancel out to $1$. So we are left with $\cos 1^\circ - \cos 3^\circ$. $(9)$

 Next we expand $(\sin 1^\circ)(2\sin 4^\circ)$. We use the same steps or strategy as we did above and get $\cos 3^\circ - \cos 5^\circ$. $(10)$

 We may be noticing a pattern. Just to make sure it is true, we will expand $(\sin 1^\circ)(2\sin 6^\circ)$. When we expand this using the same strategy above, we get $\cos 5^\circ - \cos 7^\circ$. $(11)$ So our pattern is that $(\sin 1^\circ)(2\sin \mu^\circ) = \cos (\mu-1)^\circ - \cos (\mu+1)^\circ. (12)$

 Now we notice that when we add $(9)$ and $(10)$, $-\cos 3^\circ$ and $+\cos 3^\circ$ cancel themselves out to just $\cos 1^\circ - \cos 5^\circ$. Then, when we add $(11)$ to it, $+ \cos 5^\circ$ and $- \cos 5^\circ$ also cancel themselves out. 

 Now, we need to figure out when we should stop canceling out. We can use formula $(12)$ and replace $88$ with $\mu$ since $88$ is the last one in $(8)$. Applying this formula and simplifying gives us $\cos 87^\circ - \cos 89$, so the last term would be $- \cos 89$.

 Thus we know that everything we cancel themselves out except the first and the last term. So, we are left with 
 \[(\cos 1^\circ - \cos 89^\circ) + \sin 1^\circ = \cos 1^\circ. (13)\]

 Now we notice that we have $\cos 1^\circ$ on both sides. Thus we can subtract $\cos 1^\circ$ from both sides to get $- \cos 89^\circ + \sin 1^\circ = 0$. Now we can add $\cos 89^\circ$ to both sides and get $\sin 1^\circ = \cos 89^\circ$. Now we also know that $\sin (\omega)^\circ = \cos(90-\omega)$. Thus we know that $\sin 1^\circ = \cos 89^\circ$. Hence, our proof is complete. $\square$

 


	\item (\textit{USAMO 1996}) For any nonempty set $S$ of real numbers, let $\sigma(S)$ denote the sum of the elements of $S$. Given a set $A$ of $n$ positive integers, consider the collection of all distinct sums $\sigma(S)$ as $S$ ranges over the nonempty subsets of $A$.  Prove that this collection of sums can be partitioned into $n$ classes so that in each class, the ratio of the largest sum to the smallest sum does not exceed 2.


 



\textbf{Solution 1}

Let set $A$ consist of the integers $a_1\le a_2\le a_3\le\dots\le a_n$. For $k\ge 2$, call $a_k$ greedy if $\sum_{i=1}^{k-1}a_i < a_k$. Also call $a_1$ greedy. Now put all elements of $A$ into groups of consecutive terms in such a way that each group $G$ begins with a greedy term, call it $a_p$, and ends on the term $a_{q-1}$ just before the next greedy term after $a_p$. (If $a_p$ is the last greedy term, let $q-1=n$.) We introduce some more terminology. A sum $\sigma(S)$ is said to "belong to" a group $G$ if $\max(S)\in G$. Denote by $\mathcal{S}(G)$ the set of all sums belonging to $G$.

 \textbf{We now show that we can divide $\mathcal{S}(G)$ into $|G|$ (the cardinality of $G$) classes in such a way that the maximum and minimum sum belonging to each class differ by no more than a factor of 2.} Using the previous notation, we first prove that $\frac{\max{\mathcal{S}(G)}}{\min{\mathcal{S}(G)}}< 2^{|G|}=2^{q-p}$. Note that 
 \[\max{\mathcal{S}(G)}=\sum_{i=1}^{q-1}a_i\] and 
 \[\min{\mathcal{S}(G)}=a_p.\] Taking note of that fact that $a_{p+1},a_{p+2},\dots,a_{q-1}$ are not greedy numbers, we write:

 
 \begin{align*} \max{\mathcal{S}}&=\sum_{i=1}^{q-1}a_i \\ &=a_{q-1}+\sum_{i=1}^{q-2}a_i \\ &\le 2\sum_{i=1}^{q-2}a_i\\ &\le 2^2\sum_{i=1}^{q-3}a_i\\ &\vdots\\ &\le 2^{q-p-1}\sum_{i=1}^{p}a_i\\ &= 2^{q-p-1}(\sum_{i=1}^{p-1}a_i + a_p)\\ &<2^{q-p-1}(2a_p)\\ &=2^{q-p}\min\mathcal{S} \end{align*}


 where the inequalities after the ellipses result from the fact that $a_p$ is a greedy number (which implies by definition that $\sum_{i=1}^{p-1}a_i<a_p$). This proves the desired inequality. Now we can prove the result in bold above. Divide $\mathcal{S}(G)$ into $|G|=q-p$ classes by taking those terms in $[a_p,2a_p)$ and placing them in the first class, taking those terms in $[2a_p,2^2a_p)$ and placing them in another, and so on, until we reach $[2^{q-p-1}a_p,2^{q-p}a_p)$. The inequality we proved above shows that all of the sums in $\mathcal{S}(G)$ will fall in one of these classes, as the intervals into which the classes fall form a continuous range bounded by $\min\mathcal{S}(G)=a_p$ on the bottom and $2^{q-p}a_p>\max\mathcal{S}(G)$ on the top. This proves the result in bold.

 However, that clearly implies the desired conclusion, as every sum belongs to a group, and every sum belonging to a group is a member of a class. Moreover, there will be precisely $n$ classes, because every term $a_i$ belongs to a group and for each group $G$, there are as many classes for $\mathcal{S}(G)$ as there are terms in $G$.

 



\textbf{Solution 2}

We proceed to solve this problem by induction. Our base case is when $A$ has 1 integer - there is 1 class: itself.

 Let $A$ now have $n$ elements, and assume it works for this case. We add $k$, another element, to the set. WLOG, let $k$ be the maximum element. Note that $\sigma(A)$ has to be the maximum element of a class.

 We always get a new class that works: $\sigma(A+\{k\})$ being the maximum element. This way all elements including $k$ will be in that class if $k>\sigma(A)$. When $k<\sigma(A)$ it also works since $k$ must be the maximum element from our WLOG condition above, so our induction step is done.

 


	\item (\textit{USAMO 1996}) Let $ABC$ be a triangle.  Prove that there is a line $\ell$ (in the plane of triangle $ABC$) such that the intersection of the interior of triangle $ABC$ and the interior of its reflection $A'B'C'$ in $\ell$ has area more than $\frac23$ the area of triangle $ABC$.


 



\textbf{Solution}

Let the triangle be $ABC$. Assume $A$ is the largest angle. Let $AD$ be the altitude. Assume $AB \le AC$, so that $BD \le BC/2$. If $BD > \frac{BC}{3}$, then reflect in $AD$. If $B'$ is the reflection of $B$, then $B'D = BD$ and the intersection of the two triangles is just $ABB'$. But $BB' = 2BD > \frac{2}{3} BC$, so $ABB'$ has more than $\frac{2}{3}$ the area of $ABC$.

 If $BD < BC/3$, then reflect in the angle bisector of $C$. The reflection of $A'$ is a point on the segment $BD$ and not $D$. (It lies on the line $BC$ because we are reflecting in the angle bisector. $A'C > DC$ because $\angle{CAD} < \angle{CDA} = 90^{\circ}$. Finally, $A'C \le BC$ because we assumed $\angle B$ does not exceed $\angle A$). The intersection is at least $AA'C$. But  $\frac{[AA'C]}{[ABC]} = \frac{CA'}{CB} > CD/CB \ge 2/3$.

 


	\item (\textit{USAMO 1996}) An $n$-term sequence $(x_1, x_2, \ldots, x_n)$ in which each term is either 0 or 1 is called a \textit{binary sequence of length} $n$.  Let $a_n$ be the number of binary sequences of length $n$ containing no three consecutive terms equal to 0, 1, 0 in that order.  Let $b_n$ be the number of binary sequences of length $n$ that contain no four consecutive terms equal to 0, 0, 1, 1 or 1, 1, 0, 0 in that order. Prove that $b_{n+1} = 2a_n$ for all positive integers $n$.


 



\textbf{Solution}

Given any binary sequence $B=(b_1,b_2,b_3,\dots,b_k)$, define $f(B)=(|b_2-b_1|,|b_3-b_2|,\dots,|b_k-b_{k-1}|)$. The operator $f$ basically takes pairs of consecutive terms and returns 0 if the terms are the same and 1 otherwise. Note that for every sequence $S$ of length $n$ there exist exactly two binary sequences $B$ of length $n+1$ such that $f(B)=S$.

 If $f(B)$ does not contain the string 0, 1, 0, $B$ cannot contain either of the strings 0, 0, 1, 1 or 1, 1, 0, 0. Conversely, if $B$ does not contain the sequences 0, 0, 1, 1 or 1, 1, 0, 0, $f(B)$ cannot contain 0, 1, 0. There are $a_n$ such $f(B)$ and $b_{n+1}$ such $B$. Since each $S$ corresponds with two $B$, there are twice as many such $B$ as such $S$; thus, $b_{n+1}=2a_n$.

 


	\item (\textit{USAMO 1996}) Let $ABC$ be a triangle, and $M$ an interior point such that $\angle MAB=10^\circ$, $\angle MBA=20^\circ$, $\angle MAC=40^\circ$ and $\angle MCA=30^\circ$. Prove that the triangle is isosceles.


 



\textbf{Solution 1}

Clearly, $\angle AMB = 150^\circ$ and $\angle AMC = 110^\circ$. Now by the Law of Sines on triangles $ABM$ and $ACM$, we have 
 \[\frac{AB}{\sin 150^\circ} = \frac{AM}{\sin 20^\circ}\] and 
 \[\frac{AC}{\sin 110^\circ} = \frac{AM}{\sin 30^\circ}.\] Combining these equations gives us 
 \[\frac{AB}{AC} = \frac{\sin 150^\circ \sin 30^\circ}{\sin 20^\circ \sin 110^\circ}.\] Without loss of generality, let $AB = \sin 150^\circ \sin 30^\circ = \frac{1}{4}$ and $AC = \sin 20^\circ \sin 110^\circ$. Then by the Law of Cosines, we have

 
 \begin{align*} BC^2 &= AB^2 + AC^2 - 2(AB)(BC)\cos\angle BAC\\ &= \frac{1}{16} + \sin^2 20^\circ\sin^2 110^\circ - 2\left(\frac{1}{4}\right)\sin 20^\circ\sin 110^\circ\cos 50^\circ \\ &= \frac{1}{16} + \sin^2 20^\circ \sin^2 110^\circ - \frac{1}{2}\sin 20^\circ\sin 110^\circ\sin 40^\circ \\ &= \frac{1}{16} + \sin^2 20^\circ \sin^2 110^\circ - \sin 20^\circ\sin 110^\circ\sin 20^\circ\cos 20^\circ \\ &= \frac{1}{16} \end{align*}


 Thus, $AB = BC$, our desired conclusion.

 



\textbf{Solution 2}


                    \begin{figure}[H]
                        \centering
                        \includegraphics[width=0.5\textwidth]{Solution/USAMO_I_1996_5_Solution_Figure_1.png}
                    \end{figure}
                    

 

 By the law of sines, $\frac{BM}{sin(10^\circ)}=\frac{AM}{sin(20^\circ)}$ and $\frac{CM}{sin(40^\circ)}=\frac{AM}{sin(30^\circ)}$, so $\frac{BM}{CM}=\frac{sin(10^\circ)sin(30^\circ)}{sin(20^\circ)sin(40^\circ)}$.

 Let $\angle MBC=x$. Then, $\angle MCB=80^\circ-x$. By the law of sines, $\frac{BM}{CM}=\frac{sin(80^\circ-x)}{sin(x)}$.

 Combining, we have $\frac{sin(80^\circ-x)}{sin(x)}=\frac{sin(10^\circ)sin(30^\circ)}{sin(20^\circ)sin(40^\circ)}$. From here, we can use the given trigonometric identities at each step:

 
 \begin{align*} \frac{sin(80^\circ-x)}{sin(x)}&=\frac{sin(10^\circ)sin(30^\circ)}{sin(20^\circ)sin(40^\circ)}\\ sin(80^\circ-x)sin(20^\circ)sin(40^\circ)&=sin(10^\circ)sin(30^\circ)sin(x)\\ sin(80^\circ-x)sin(20^\circ)sin(40^\circ)&=\frac{1}{2}sin(10^\circ)sin(x)&[sin(30^\circ)=1/2]\\ sin(80^\circ-x)sin(30^\circ-10^\circ)sin(30^\circ+10^\circ)&=\frac{1}{2}sin(10^\circ)sin(x)\\ sin(80^\circ-x)(cos^2(10^\circ)-cos^2(30^\circ))&=\frac{1}{2}sin(10^\circ)sin(x)&[sin(A-B)sin(A+B)=cos^2 B-cos^2 A]\\ sin(80^\circ-x)(cos^2(10^\circ)-\frac{3}{4})&=\frac{1}{2}sin(10^\circ)sin(x)&[cos(30^\circ)=\frac{\sqrt{3}}{2}]\\ sin(80^\circ-x) \frac{4cos^3(10^\circ)-3cos(10^\circ)}{4cos(10^\circ)}&=\frac{1}{2}sin(10^\circ)sin(x)\\ sin(80^\circ-x) \frac{cos(30^\circ)}{4cos(10^\circ)}&=\frac{1}{2}sin(10^\circ)sin(x)&[cos(3A)=4cos^3 A-3cos A]\\ sin(80^\circ-x)cos(30^\circ)&=2sin(10^\circ)cos(10^\circ)sin(x)\\ sin(80^\circ-x)cos(30^\circ)&=sin(20^\circ)sin(x)&[sin(2A)=2sin A cos A ]\\ sin(80^\circ-x)sin(60^\circ)&=sin(20^\circ)sin(x)&[cos(30^\circ)=sin(60^\circ)]\\ \frac{1}{2}(cos(20^\circ-x)-cos(140^\circ-x))&=\frac{1}{2}(cos(20^\circ-x)-cos(20^\circ+x))&[sin A sin B=\frac{1}{2}(cos(A-B)-cos(A+B))]\\ cos(140^\circ-x)&=cos(20^\circ+x) \end{align*}


 The only acute angle satisfying this equality is $x=60^\circ$. Therefore, $\angle ACB=80^\circ-x+30^\circ=50^\circ$ and $\angle BAC=10^\circ+40^\circ=50^\circ$. Thus, $\triangle ABC$ is isosceles.

 



\textbf{Solution 3}

If $\angle{MBC} = x$ then by Angle Sum in a Triangle we have $\angle{MCB} = 80^\circ - x$. By Trig Ceva we have
 \[\sin 10^\circ \sin x \sin 30^\circ = \sin (80^\circ - x) \sin 40^\circ \sin 20^\circ.\]Because $\dfrac{\sin x}{\sin (80^\circ - x)}$ is monotonic increasing over $(0, \dfrac{\pi}{2})$, there is only one solution $0 \le x \le \dfrac{\pi}{2}$ to the equation. We claim it is $x = 60^\circ$, which will make $ABC$ isosceles with $\angle{A} = \angle{C}$.

 Notice that
 \[\sin 20^\circ \sin 20^\circ \sin 40^\circ = 2 \sin 10^\circ \cos 10^\circ \sin 20^\circ \sin 40^\circ\]
 \[= \sin 10^\circ (\sin 10^\circ + \frac{1}{2}) \sin 40^\circ\]
 \[= \sin 10^\circ (\frac{1}{2} \sin 40^\circ + \frac{1}{2} (\cos 30^\circ - \cos 50^\circ))\]
 \[= \sin 10^\circ \frac{1}{2} \cos 30^\circ\]
 \[= \sin 10^\circ \sin 30^\circ \sin 60^\circ,\]as desired.

 


	\item (\textit{USAMO 1996}) Determine (with proof) whether there is a subset $X$ of the integers with the following property: for any integer $n$ there is exactly one solution of $a + 2b = n$ with $a,b \in X$.


 



\textbf{Solution}

Start with an incomplete subset $S = (S_1, S_2, S_3, ... S_m)$, such that for any integer n, there is exactly zero or one solutions to $a + 2b = n$ with $a,b \in S$. Let $N$ be the smallest integer such that for any $S_i$, $|S_i| < N$. Note that $|S_i+2S_j| < 3N$ for any $S_i$ and $S_j$

 
Suppose $M$ is the smallest non-negative integer without a solution in $S$ yet. Clearly, $0 \le M \le 3N$. Generate $S_{m+1}$ and $S_{m+2}$ such that $S_{m+1} = -10N - M$, and $S_{m+2} = 5N + M$. Thus, we now have the solution $S_{m+1}+2S_{m+2} = M$.

 \textit{Note: The values 10 and 5 can be replaced by any sufficiently large values such that the first is twice the second.}

 
Now, we must prove that the addition of these two terms to $S$ does not result in an integer n that has two solutions. Of course, $S_{m+1} + 2S_{m+2} = M$ which previously had no solutions. Furthermore, $S_{m+1} + 2S_{m+2} = -15N - M$.

 
            \begin{enumerate}[label=(\roman*)]
                \item For any $S_i$, $-12N - M < S_{m+1}+2S_i < -8N-M$, and $-21N - 2M < S_i+2S_{m+1} < -19N - 2M$.
            \end{enumerate}
            
            \begin{enumerate}[label=(\roman*)]
                <dd>
            \begin{enumerate}[label=(\roman*)]
                \item Since $0 \le M \le 3N$, we get that  $-18N < S_{m+1}+2S_i < -8N$ and $-27N < S_i+2S_{m+1} < -19N$
            \end{enumerate}
            </dd>
            \end{enumerate}
            
            \begin{enumerate}[label=(\roman*)]
                \item Similarly, $3N + M < S_{m+2}+2S_i < 5N + M$, and $9N+2M < S_i+2S_{m+2} < 11N+2M$.
            \end{enumerate}
            
            \begin{enumerate}[label=(\roman*)]
                <dd>
            \begin{enumerate}[label=(\roman*)]
                \item Since $0 \le M \le 3N$, we get that  $3N < S_{m+2}+2S_i < 8N$ and $9N < S_i+2S_{m+2} < 17N$.
            \end{enumerate}
            </dd>
            \end{enumerate}
            Since all of these sums (other than $M$) are either greater than $3N$ or less than $-3N$, they are all sums that previously had no solutions. Furthermore, none of these sums are duplicated, as sums of different forms are contained in disjoint ranges of integers.

 Thus, we have proved that we can generate a subset $S$ such that all non-negative integers n have a unique solution $a + 2b = n$.

 
For negative integers M that have no solutions in $S$ a similar proof holds, but instead generating the terms $S_{m+1} = 10N - M$ and $S_{m+2} = -5N + M$.

 
For any integer M that currently has no solution in S, we can always add two terms $S_{m+1}$ and $S_{m+2}$ such that $S_{m+1} + 2S_{m+2} = M$ that do not result in duplicated sums.

 Thus, there does exist a subset $X$ of the integers such that for any integer $n$ there is exactly one solution to $a + 2b = n$ with $a, b \in X$.

 


\end{enumerate}